%! Transformations of Functions — Unit 1, Lesson 1.4 — Homework
\documentclass[10pt]{article}
\usepackage{precalculus-article}
\usepackage{precalculus-boxes}
\newcommand{\TallMath}[1]{$\displaystyle #1\rule[-1.4em]{0pt}{3.2em}$}

\begin{document}

\pageheader{Unit 1, Lesson 1.4}{Homework}
\namedateperiod

\vspace{0.12in}

\begin{enumerate}[itemsep=18pt]

\item For each transformation rule below, identify the type of transformation (vertical
      translation, horizontal translation, vertical dilation, or horizontal dilation/reflection)
      and describe its effect on the graph of $f$.

      \begin{enumerate}[label=(\alph*), itemsep=10pt]
        \item $g(x) = f(x) - 5$

              Type: \blank{4cm} \quad Effect: \blank{6cm}

        \item $g(x) = f(x - 3)$

              Type: \blank{4cm} \quad Effect: \blank{6cm}

        \item $g(x) = \dfrac{1}{3}f(x)$

              Type: \blank{4cm} \quad Effect: \blank{6cm}

        \item $g(x) = f(-x)$

              Type: \blank{4cm} \quad Effect: \blank{6cm}
      \end{enumerate}

\item The table below gives values of $f$.

      \medskip
      \renewcommand{\arraystretch}{1.3}
      \begin{tabular}{|c|c|c|c|c|c|}
      \hline
      $x$    & $-2$ & $-1$ & $0$ & $1$ & $2$ \\
      \hline
      $f(x)$ & $8$  & $2$  & $0$ & $2$ & $8$ \\
      \hline
      \end{tabular}

      \medskip
      Complete a table of values for each transformed function.

      \medskip
      \textbf{(a)} $g(x) = f(x) + 4$

      \smallskip
      \renewcommand{\arraystretch}{1.3}
      \begin{tabular}{|c|c|c|c|c|c|}
      \hline
      $x$    & $-2$ & $-1$ & $0$ & $1$ & $2$ \\
      \hline
      $g(x)$ & \blank{1.2cm} & \blank{1.2cm} & \blank{1.2cm} & \blank{1.2cm} & \blank{1.2cm} \\
      \hline
      \end{tabular}

      \medskip
      \textbf{(b)} $h(x) = -f(x)$

      \smallskip
      \renewcommand{\arraystretch}{1.3}
      \begin{tabular}{|c|c|c|c|c|c|}
      \hline
      $x$    & $-2$ & $-1$ & $0$ & $1$ & $2$ \\
      \hline
      $h(x)$ & \blank{1.2cm} & \blank{1.2cm} & \blank{1.2cm} & \blank{1.2cm} & \blank{1.2cm} \\
      \hline
      \end{tabular}

\item Let $f$ have domain $[1,\,5]$ and range $[-3,\,6]$. For each transformed function,
      state the domain and range.

      \begin{enumerate}[label=(\alph*), itemsep=10pt]
        \item $g(x) = f(x) + 2$ \quad Domain: \blank{4cm} Range: \blank{4cm}
        \item $g(x) = f(x + 4)$ \quad Domain: \blank{4cm} Range: \blank{4cm}
        \item $g(x) = 3f(x)$    \quad Domain: \blank{4cm} Range: \blank{4cm}
        \item $g(x) = f(x-1)-2$ \quad Domain: \blank{4cm} Range: \blank{4cm}
      \end{enumerate}

\item Describe all transformations applied to $f$ to obtain $g(x) = -3f(x+2) + 1$.
      List them in order of application.

      \vspace{0.08in}
      \writelines{3}

\item \textbf{(AP-Style Multiple Choice)} Let $f$ be a function with domain $[-4,\,4]$ and
      range $[0,\,8]$. Which of the following gives the range of $g(x) = -2f(x) + 1$?

      \medskip
      \begin{enumerate}[label=(\Alph*), itemsep=4pt]
        \item $[-15,\,1]$
        \item $[-16,\,0]$
        \item $[1,\,17]$
        \item $[-15,\,-1]$
      \end{enumerate}

      \vspace{0.06in}
      Answer: \blank{1.5cm} \quad Show your reasoning:

      \vspace{0.08in}
      \writelines{2}

\end{enumerate}

\begin{reflectionbox}
What is the trickiest part of horizontal transformations? Why does $f(x+h)$ shift
\textit{opposite} to the sign of $h$?

\vspace{0.06in}
\writelines{2}
\end{reflectionbox}

\end{document}
